\label{sec:outline}

The remainder of this dissertation is organized as follows.

Chapter 2 introduces novelty in open-world autonomy. It contrasts closed-world assumptions with open-world settings, develops the view of novelty as an agent-relative mismatch, presents the novelty taxonomy used throughout the dissertation, and introduces assimilation and accommodation as two modes of response to mismatch. The chapter closes by identifying the requirements that open-world agents must satisfy: detecting and characterizing mismatch, adapting from limited experience, revising representations when needed, and integrating perception, reasoning, learning, and action.

Chapter 3 reviews related work on adaptation under novelty. It organizes existing approaches by whether they adapt within a fixed representation or extend the agent's representational structure, and considers foundation models as a source of broad priors and candidate structure. The chapter identifies the gap addressed by this dissertation: structural adaptation and generalization that are localizable, grounded in action, and mediated by explicit abstractions.

Chapter 4 presents the formal preliminaries and structured abstraction framework used in the dissertation. It defines the agent--environment formalism, symbolic planning, reinforcement learning, hybrid planning--learning, and the operator--executor structure used in the technical chapters. It then develops the role of structured abstractions in mediating the agent--environment interaction, establishing the formal vocabulary that the subsequent chapters build upon.

Chapter 5 introduces benchmark environments, formalism, and evaluation methodologies for systematically studying novelty. These environments enable controlled injection of diverse classes of mismatch and provide metrics for analyzing how agent performance degrades and recovers under changing conditions. This chapter establishes an experimental foundation for evaluating adaptation and generalization in open-world settings.

Chapter 6 investigates adaptation under novelty. It presents a neurosymbolic cognitive architecture for open-world novelty handling and hybrid planning--learning methods that enable agents to recover from execution failures by leveraging structured action abstractions. These approaches utilize counterfactual reasoning based methods to efficiently identify the source of failure and acquire missing behaviors, enabling agents to adapt to novel conditions without relearning entire tasks. By localizing failures to specific components of a plan and enabling targeted learning, these methods allow agents to efficiently acquire missing behaviors without relearning entire tasks.

Chapter 7 introduces FLEX, an object-centric force-space framework for articulated object manipulation. FLEX shows how policies trained in force space can generalize across articulated objects and transfer directly to different robotic platforms without retraining.

Chapter 8 extends the force-space interaction abstraction to rigid body transport. It presents object-centric force policies for pushing heavy objects under unknown mass and friction, demonstrating zero-shot sim-to-real transfer on a legged mobile manipulator.

In Chapters 5 through 8, we also review related work specific to each technical contribution and discuss how the proposed methods compare to and complement existing approaches.

Chapter 9 concludes the dissertation by summarizing the main findings, articulating a unified perspective on how structured abstractions mediate an agent's response to novelty, and discussing future directions for research on autonomous agents and robots operating in open-world environments.
